\chapter*{Definición de Huffman}
\addcontentsline{toc}{chapter}{Definición de Huffman}

\noindent
El algoritmo de Huffman es un método de \textbf{compresión sin pérdida} que asigna códigos binarios de longitud variable a los símbolos de una fuente, otorgando códigos más cortos a los símbolos más frecuentes y códigos más largos a los menos frecuentes. El resultado es un \textbf{código prefijo}: ningún código es prefijo de otro, lo que garantiza una decodificación unívoca sin necesidad de delimitadores.

\section*{Procedimiento del algoritmo}
\addcontentsline{toc}{section}{Procedimiento del algoritmo}

\noindent
El algoritmo opera de la siguiente manera:
\begin{enumerate}
    \item Se calcula la frecuencia de aparición de cada símbolo en el archivo.
    \item Se construye un \textit{min-heap} (cola de prioridad) donde cada símbolo es un nodo hoja con su frecuencia asociada.
    \item Repetidamente, se extraen los dos nodos con menor frecuencia, se combinan en un nuevo nodo padre cuya frecuencia es la suma de ambos, y se reinserta en el heap.
    \item El proceso se repite hasta que queda un único nodo raíz: el \textbf{árbol de Huffman}.
    \item Se recorre el árbol asignando ``0'' a cada rama izquierda y ``1'' a cada rama derecha, generando así el código binario de cada símbolo.
\end{enumerate}

\section*{Longitud promedio de palabra código}
\addcontentsline{toc}{section}{Longitud promedio de palabra código}

\noindent
La eficiencia de la compresión se mide mediante la longitud promedio de palabra código:
\begin{equation}
L_{av} = \sum_{i=1}^{q} p_i \cdot l_i
\end{equation}
donde $p_i$ es la probabilidad del símbolo $i$ y $l_i$ la longitud de su palabra código. Huffman produce la asignación óptima que minimiza $L_{av}$, alcanzando valores cercanos a la entropía de la fuente.
